```latex
\documentclass[11pt]{article}
\usepackage{amsmath,amssymb,amsthm,mathrsfs}

\newtheorem{theorem}{Theorem}[section]
\newtheorem{lemma}[theorem]{Lemma}
\newtheorem{proposition}[theorem]{Proposition}
\newtheorem{corollary}[theorem]{Corollary}
\theoremstyle{definition}
\newtheorem{definition}[theorem]{Definition}
\theoremstyle{remark}
\newtheorem{remark}[theorem]{Remark}

\newcommand{\Lie}{\mathcal{L}}
\newcommand{\D}{\mathscr{D}}
\newcommand{\R}{\mathbb{R}}
\newcommand{\Z}{\mathbb{Z}}
\newcommand{\supp}{\operatorname{supp}}
\newcommand{\htop}{h_{\mathrm{top}}}

\begin{document}

\title{Globally hypoelliptic vector fields have zero topological entropy}
\author{}
\date{}
\maketitle

\begin{abstract}
Let $X$ be a $C^\infty$ vector field on a compact connected smooth manifold $M$ and $g_t$ its flow.
Assume that the cohomological operator $\Lie_X$ is globally hypoelliptic on scalar distributions: whenever $u\in\D'(M)$ satisfies $\Lie_X u=f$ with $f\in C^\infty(M)$, then $u$ is a smooth function.
We prove that the flow $(g_t)$ has zero topological entropy.
The proof has two ingredients: global hypoellipticity excludes periodic orbits in dimension $\ge 2$, and a theorem of Katok implies that any $C^{1+\alpha}$ diffeomorphism with positive entropy must have periodic points.
\end{abstract}

\section{Setting}

Let $M$ be a compact connected $C^\infty$ manifold of dimension $n:=\dim M$, and let
$X\in\mathfrak{X}^\infty(M)$ be a smooth vector field.

\subsection{Completeness of the flow}

\begin{lemma}[Completeness]\label{lem:complete}
The vector field $X$ is complete: it generates a global $C^\infty$ flow $(g_t)_{t\in\R}$.
\end{lemma}

\begin{proof}
By the standard local existence and uniqueness theorem for ODEs on manifolds, through each $p\in M$
there is a unique maximal integral curve $\gamma_p:(a_p,b_p)\to M$ with $\gamma_p(0)=p$ and
$\dot{\gamma}_p(t)=X(\gamma_p(t))$.

We show $b_p=+\infty$; the argument for $a_p=-\infty$ is analogous.
Assume for contradiction that $b_p<+\infty$.
Choose a sequence $t_k\uparrow b_p$.
By compactness of $M$, the sequence $\gamma_p(t_k)$ has a convergent subsequence
(still denoted $\gamma_p(t_k)$) with $\gamma_p(t_k)\to q\in M$.

Fix a chart $(U,\psi)$ about $q$.
For $k$ large, $\gamma_p(t_k)\in U$.
In local coordinates, the ODE becomes $\dot{x}=F(x)$ where $F$ is smooth on the open set $\psi(U)\subset\R^n$.
Choose a compact set $K\Subset \psi(U)$ containing $\psi(q)$.
For $k$ large enough, $\psi(\gamma_p(t_k))\in K$.
Since $F$ is smooth, it is Lipschitz on $K$, hence solutions starting in $K$ exist on a uniform time interval
$[-\varepsilon,\varepsilon]$ (depending only on $K$ and $F$).
Thus the integral curve through $\gamma_p(t_k)$ extends at least up to time $t_k+\varepsilon$.
Choosing $k$ so large that $t_k+\varepsilon>b_p$ contradicts maximality of $\gamma_p$.
Hence $b_p=+\infty$.
\end{proof}

We denote by $(g_t)_{t\in\R}$ the resulting global flow.

\subsection{Distributions and Lie derivative}

Let $\D'(M)$ denote the space of scalar distributions on $M$, i.e.\ the continuous dual of $C^\infty(M)$
(with its usual Fr\'echet topology).
Fix a smooth positive density $\Omega$ on $M$.
We identify $C^\infty(M)$ with a subspace of $\D'(M)$ via
\[
f\in C^\infty(M)\ \longmapsto\ f\Omega\in\D'(M),\qquad
(f\Omega)(\varphi):=\int_M f\varphi\,\Omega,\quad \varphi\in C^\infty(M).
\]
A distribution $u$ is called \emph{smooth} if $u=f\Omega$ for some $f\in C^\infty(M)$.

The Lie derivative $\Lie_X$ acts on test functions by $\Lie_X\varphi=X(\varphi)$.
On distributions it is defined by duality:
\begin{equation}\label{eq:LieDist}
(\Lie_X u)(\varphi):=-u(\Lie_X\varphi),\qquad u\in\D'(M),\ \varphi\in C^\infty(M).
\end{equation}

\begin{definition}[Global hypoellipticity]\label{def:GH}
We say that $X$ is \emph{globally hypoelliptic} if for every $u\in\D'(M)$ and $f\in C^\infty(M)$,
\[
\Lie_X u = f\Omega \quad\Longrightarrow\quad u \text{ is smooth}.
\]
\end{definition}

This is exactly the hypothesis in the problem statement (after the identification of $f$ with $f\Omega$).

\subsection{Topological entropy}
We use the Bowen--Dinaburg definition of entropy for continuous maps on compact metric spaces.
If $F:(Y,d)\to(Y,d)$ is continuous and $n\ge 1$, define
\[
d_n(x,y):=\max_{0\le k<n} d(F^k x, F^k y).
\]
A set $E\subset Y$ is \emph{$(n,\varepsilon)$-separated} if $d_n(x,y)\ge\varepsilon$ for all distinct $x,y\in E$.
Let $s_n(\varepsilon)$ be the maximum cardinality of an $(n,\varepsilon)$-separated set.
Then
\[
\htop(F):=\lim_{\varepsilon\downarrow 0}\ \limsup_{n\to\infty}\ \frac1n\log s_n(\varepsilon).
\]

For a continuous flow $(g_t)_{t\in\R}$ on a compact metric space, one sets $\htop(g_t):=\htop(g_1)$.
This agrees with the intrinsic definition of entropy for flows and satisfies the scaling law
$\htop(g_t)=|t|\htop(g_1)$; see Bowen--Walters \cite{BW}.
Thus it suffices to prove $\htop(g_1)=0$.

\section{Two support lemmas}

\begin{lemma}\label{lem:smoothsupport}
Let $f\in C^\infty(M)$ with $f\not\equiv 0$.
Then the support of the distribution $f\Omega$ has nonempty interior.
\end{lemma}

\begin{proof}
Choose $x_0\in M$ with $f(x_0)\neq 0$.
By continuity, there exists an open neighborhood $U$ of $x_0$ such that $f|_U$ is nowhere zero.
Pick any $\varphi\in C^\infty(M)$ supported in $U$ with $\varphi\ge 0$ and $\varphi\not\equiv 0$.
Then $\int_M f\varphi\,\Omega=\int_U f\varphi\,\Omega\neq 0$, so $f\Omega$ does not vanish on $U$.
Hence $U\subset\supp(f\Omega)$.
\end{proof}

\begin{lemma}\label{lem:submanifoldempty}
If $N\subset M$ is an embedded $C^1$ submanifold with $\dim N<\dim M$, then $N$ has empty interior in $M$.
\end{lemma}

\begin{proof}
Let $x\in N$.
By the definition of embedded submanifold there exists a chart $(U,\psi)$ about $x$ such that $\psi(U)\subset\R^n$ is open and
\[
\psi(N\cap U)=\psi(U)\cap (\R^k\times\{0\})
\]
for some $k<n$.
The set $\R^k\times\{0\}$ has empty interior in $\R^n$, hence $N\cap U$ has empty interior in $U$.
Since $x$ was arbitrary, $N$ has empty interior in $M$.
\end{proof}

\section{Periodic orbit distributions}

\subsection{Periodic orbit measures}

Suppose that for some $p\in M$ and $T>0$ one has $g_T(p)=p$.
Let $\gamma:=\{g_t(p):t\in\R\}$ be the corresponding orbit.

Define the distribution $\mu_\gamma\in\D'(M)$ by the orbit average
\begin{equation}\label{eq:orbitmeasure}
\mu_\gamma(\varphi):=\frac1T\int_0^T \varphi(g_t(p))\,dt,\qquad \varphi\in C^\infty(M).
\end{equation}
This is an invariant probability measure supported on $\gamma$.

\begin{lemma}\label{lem:orbit_Lie_zero}
One has $\Lie_X \mu_\gamma=0$.
\end{lemma}

\begin{proof}
For $\varphi\in C^\infty(M)$,
\[
(\Lie_X\mu_\gamma)(\varphi)=-\mu_\gamma(X\varphi)
=-\frac1T\int_0^T X\varphi(g_t(p))\,dt.
\]
But $\frac{d}{dt}\bigl(\varphi(g_t(p))\bigr)=X\varphi(g_t(p))$.
Therefore
\[
(\Lie_X\mu_\gamma)(\varphi)
=-\frac1T\int_0^T \frac{d}{dt}\bigl(\varphi(g_t(p))\bigr)\,dt
=-\frac1T\bigl(\varphi(g_T(p))-\varphi(p)\bigr)=0.
\]
\end{proof}

\begin{lemma}\label{lem:orbit_support}
$\supp(\mu_\gamma)=\gamma$.
\end{lemma}

\begin{proof}
If $U$ is open and $U\cap\gamma=\varnothing$, then any test function $\varphi$ supported in $U$ vanishes on $\gamma$,
so $\mu_\gamma(\varphi)=0$.
Hence $\supp(\mu_\gamma)\subset\gamma$.

Conversely, let $x\in\gamma$ and let $U$ be an open neighborhood of $x$.
Choose $t_0\in[0,T)$ with $x=g_{t_0}(p)$.
By continuity there exists $\delta>0$ such that $g_t(p)\in U$ for all $t$ in a small interval around $t_0$ (mod $T$).
Pick $\varphi\in C^\infty(M)$ with $\varphi\ge 0$, $\supp(\varphi)\subset U$, and $\varphi(x)>0$.
Then $\mu_\gamma(\varphi)>0$, so $x\in\supp(\mu_\gamma)$.
\end{proof}

\subsection{No periodic orbits in dimension $\ge 2$}

\begin{proposition}\label{prop:no_periodic_dimge2}
Assume $\dim M\ge 2$ and $X$ is globally hypoelliptic in the sense of Definition~\ref{def:GH}.
Then the flow $(g_t)$ has no periodic orbits (and in particular no fixed points).
\end{proposition}

\begin{proof}
Assume for contradiction that $g_T(p)=p$ for some $p\in M$ and $T>0$, and let $\gamma$ be the corresponding orbit.
By Lemma~\ref{lem:orbit_Lie_zero}, $\Lie_X\mu_\gamma=0$.
Since $0$ is smooth, global hypoellipticity forces $\mu_\gamma$ to be smooth, i.e.\ $\mu_\gamma=f\Omega$ for some $f\in C^\infty(M)$.

Now $\mu_\gamma\neq 0$, hence $f\not\equiv 0$.
Lemma~\ref{lem:smoothsupport} implies that $\supp(\mu_\gamma)$ has nonempty interior.
But by Lemma~\ref{lem:orbit_support} one has $\supp(\mu_\gamma)=\gamma$.

The set $\gamma$ is the image of a $C^\infty$ map $t\mapsto g_t(p)$ from $\R$ to $M$ with everywhere rank $\le 1$.
If $p$ is fixed then $\gamma$ is a $0$-dimensional embedded submanifold; otherwise $\gamma$ is a periodic orbit and hence an embedded $1$-dimensional submanifold (a closed integral curve).
In either case $\dim\gamma<\dim M$ (since $\dim M\ge 2$), so Lemma~\ref{lem:submanifoldempty} gives that $\gamma$ has empty interior.
This contradicts the previous paragraph.
\end{proof}

\section{The one-dimensional case}

When $\dim M=1$, a periodic orbit may coincide with the whole manifold and the previous argument does not apply.
We treat this case separately.

\begin{proposition}\label{prop:dim1}
Assume $\dim M=1$ and $X$ is globally hypoelliptic.
Then $M\cong S^1$ and the flow $(g_t)$ is $C^\infty$ conjugate to a rigid rotation flow on $S^1$.
In particular $\htop(g_t)=0$ for all $t\in\R$.
\end{proposition}

\begin{proof}
A compact connected $1$-manifold without boundary is diffeomorphic to $S^1$ \cite[Thm.~8.1]{LeeSM}; we identify $M$ with $S^1=\R/\Z$.

\emph{Step 1: $X$ has no zeros.}
If $x_0\in S^1$ satisfies $X(x_0)=0$, then $x_0$ is a fixed point and the Dirac distribution $\delta_{x_0}$ satisfies
$\Lie_X\delta_{x_0}=0$ because
\[
(\Lie_X\delta_{x_0})(\varphi)=-\delta_{x_0}(X\varphi)=-(X\varphi)(x_0)=0.
\]
Global hypoellipticity would force $\delta_{x_0}$ to be smooth, contradicting Lemma~\ref{lem:smoothsupport} since $\supp(\delta_{x_0})=\{x_0\}$ has empty interior.
Thus $X$ is nowhere vanishing.

After possibly reversing orientation, write $X=a(\theta)\,\partial_\theta$ with $a\in C^\infty(S^1)$ strictly positive.

\emph{Step 2: Conjugacy to constant speed rotation.}
Define
\[
T:=\int_0^1 \frac{1}{a(s)}\,ds>0,\qquad
H(\theta):=\frac{1}{T}\int_0^\theta \frac{1}{a(s)}\,ds\quad(\theta\in\R).
\]
Because $a$ is $1$-periodic, $H(\theta+1)=H(\theta)+1$, so $H$ descends to a $C^\infty$ diffeomorphism $h:S^1\to S^1$.
Since $H'(\theta)=\frac{1}{Ta(\theta)}>0$, $H$ is strictly increasing.
A direct computation yields
\[
h_*X = a(\theta)H'(\theta)\,\partial_\vartheta = \frac{1}{T}\,\partial_\vartheta,
\qquad \vartheta=h(\theta).
\]
Therefore $h\circ g_t = R_t\circ h$ for all $t$, where $R_t(\vartheta)=\vartheta+t/T$ is the rigid rotation flow.

\emph{Step 3: Entropy.}
The time--one map $R_1$ is an isometry of $S^1$ for the standard metric, hence has zero topological entropy:
for an isometry $F$ one has $d_n=d$ for all $n$, so $s_n(\varepsilon)$ does not depend on $n$ and $\htop(F)=0$.
Entropy is invariant under topological conjugacy (see e.g.\ \cite[\S7]{Walters}), hence $\htop(g_1)=\htop(R_1)=0$.
The scaling law for flows \cite{BW} then gives $\htop(g_t)=0$ for all $t$.
\end{proof}

\section{Positive entropy forces periodic points}

We now recall two standard results.

\begin{lemma}\label{lem:positive_entropy_measure}
Let $F$ be a continuous map of a compact metric space.
If $\htop(F)>0$, then there exists an ergodic $F$-invariant Borel probability measure $\mu$ with $h_\mu(F)>0$.
\end{lemma}

\begin{proof}[Citation]
By the variational principle (Walters \cite[\S8.2]{Walters}), there exists an $F$-invariant Borel probability measure $\nu$
with $h_\nu(F)>0$.
By the ergodic decomposition theorem and affinity of metric entropy (Walters \cite[\S8.3]{Walters}),
some ergodic component $\mu$ of $\nu$ satisfies $h_\mu(F)>0$.
\end{proof}

\begin{theorem}[Katok]\label{thm:katok}
Let $F$ be a $C^{1+\alpha}$ diffeomorphism of a compact smooth manifold, $\alpha>0$.
If there exists an ergodic $F$-invariant probability measure $\mu$ with $h_\mu(F)>0$, then $F$ has a periodic point.
\end{theorem}

\begin{proof}[Citation]
This follows from Katok's theorem that positive-entropy measures can be approximated by hyperbolic horseshoes,
and horseshoes contain periodic points; see Katok \cite{Katok80}.
\end{proof}

\begin{corollary}\label{cor:positive_entropy_periodic}
Let $F$ be a $C^{1+\alpha}$ diffeomorphism of a compact manifold.
If $\htop(F)>0$, then $F$ has a periodic point.
\end{corollary}

\begin{proof}
If $\htop(F)>0$, Lemma~\ref{lem:positive_entropy_measure} produces an ergodic invariant measure with positive metric entropy,
and Theorem~\ref{thm:katok} yields a periodic point.
\end{proof}

\section{Main theorem}

\begin{theorem}\label{thm:main}
Let $M$ be a compact connected smooth manifold and $X$ a smooth vector field on $M$.
Assume that $X$ is globally hypoelliptic in the sense of Definition~\ref{def:GH}.
Then the flow $(g_t)$ generated by $X$ has zero topological entropy:
\[
\htop(g_t)=0\qquad\text{for all }t\in\R.
\]
\end{theorem}

\begin{proof}
By Bowen--Walters \cite{BW}, it suffices to show $\htop(g_1)=0$.

If $\dim M=1$, Proposition~\ref{prop:dim1} gives $\htop(g_1)=0$.

Assume $\dim M\ge 2$.
By Proposition~\ref{prop:no_periodic_dimge2}, the flow $(g_t)$ has no periodic orbits.
In particular, the time--one map $g_1$ has no periodic points: if $g_1^k(p)=p$ for some $k\ge 1$, then $g_k(p)=p$, so the orbit of $p$ under the flow is periodic.

If $\htop(g_1)>0$, then Corollary~\ref{cor:positive_entropy_periodic} (applied to the $C^\infty$ diffeomorphism $g_1$)
would imply that $g_1$ has a periodic point, contradicting the previous paragraph.
Therefore $\htop(g_1)=0$, hence $\htop(g_t)=0$ for all $t$.
\end{proof}

\begin{thebibliography}{99}

\bibitem{BW}
R.~Bowen and P.~Walters,
\newblock \emph{Expansive one-parameter flows},
\newblock J.\ Differential Equations \textbf{12} (1972), 180--193.

\bibitem{Katok80}
A.~Katok,
\newblock \emph{Lyapunov exponents, entropy and periodic orbits for diffeomorphisms},
\newblock Publ.\ Math.\ Inst.\ Hautes \'Etudes Sci.\ \textbf{51} (1980), 137--173.

\bibitem{LeeSM}
J.~M.~Lee,
\newblock \emph{Introduction to Smooth Manifolds},
\newblock 2nd ed., Graduate Texts in Mathematics, vol.~218, Springer, 2013.

\bibitem{Walters}
P.~Walters,
\newblock \emph{An Introduction to Ergodic Theory},
\newblock Graduate Texts in Mathematics, vol.~79, Springer, 1982.

\end{thebibliography}

\end{document}
```